\renewcommand{\chapterprefix}{Appendix}\renewcommand{\thechapter}{T}
\chapter{Solutions to the stretch exercises of Chapters 3--5}
\chaptercontents{A & Chapter 3 \\ B & Chapter 4 \\ C & Chapter 5}{Compare your own attempt line by line: did you introduce every variable, unfold every definition, and say which strategy you were using?}

\hsection{Chapter 3}
\sol{S3.1} ($\subseteq$) Let $b\in f[f^{-1}[V]]$; fix $a\in f^{-1}[V]$ with $f(a)=b$. Then $b=f(a)\in V$ and $b\in f[X]$. ($\supseteq$) Let $b\in V\cap f[X]$; fix $a\in X$ with $f(a)=b$. Since $f(a)\in V$, $a\in f^{-1}[V]$, so $b\in f[f^{-1}[V]]$.

\sol{S3.2} $(f^{-1}\circ g^{-1})\circ(g\circ f)=f^{-1}\circ(g^{-1}\circ g)\circ f=f^{-1}\circ\id_Y\circ f=\id_X$ and $(g\circ f)\circ(f^{-1}\circ g^{-1})=g\circ(f\circ f^{-1})\circ g^{-1}=\id_Z$. So $f^{-1}\circ g^{-1}$ is an inverse of $g\circ f$, and inverses are unique.

\sol{S3.3} ($\Rightarrow$) Let $f$ be injective, $U\subseteq X$, and $b\in f[X\setminus U]$: $b=f(a)$ with $a\notin U$. If $b\in f[U]$ then $b=f(u)$ with $u\in U$, and injectivity gives $a=u\in U$, a contradiction. So $b\in Y\setminus f[U]$. ($\Leftarrow$) Suppose the inclusion holds for every $U$. Let $f(a)=f(a')$ and suppose $a\ne a'$. Take $U=\{a'\}$. Then $a\in X\setminus U$, so $f(a)\in f[X\setminus U]\subseteq Y\setminus f[U]=Y\setminus\{f(a')\}$, so $f(a)\ne f(a')$: contradiction. Hence $a=a'$.

\sol{S3.4} If $n$ is even, $f(n)=n+1$ is odd, so $f(f(n))=(n+1)-1=n$. If $n$ is odd, $f(n)=n-1$ is even, so $f(f(n))=(n-1)+1=n$. Thus $f\circ f=\id_\Z$, so $f$ is its own two-sided inverse, hence a bijection.

\sol{S3.5} $g=g\circ\id_Y=g\circ(f\circ h)=(g\circ f)\circ h=\id_X\circ h=h$.

\sol{S3.6} ($\Rightarrow$) Suppose $f$ surjective and $g\circ f=h\circ f$. Let $b\in Y$; fix $a$ with $f(a)=b$. Then $g(b)=g(f(a))=h(f(a))=h(b)$. So $g=h$. ($\Leftarrow$) Contrapositive. Suppose $f$ is not surjective; fix $b\in Y\setminus f[X]$. Define $g,h:Y\to\{0,1\}$ by $g(y)=0$ for all $y$, and $h(b)=1$, $h(y)=0$ for $y\ne b$. For every $a\in X$, $f(a)\ne b$, so $h(f(a))=0=g(f(a))$; thus $g\circ f=h\circ f$. But $g(b)=0\ne1=h(b)$, so $g\ne h$. Cancellation fails.

\sol{S3.7} Let $f:X\to\PS(X)$ and $D=\{a\in X\mid a\notin f(a)\}\in\PS(X)$. Suppose $f$ is surjective; fix $d\in X$ with $f(d)=D$. Then $d\in D\Leftrightarrow d\notin f(d)\Leftrightarrow d\notin D$, which is impossible. So $f$ is not surjective.

\sol{S3.8} Let $g(x)=a+(b-a)x$. If $0<x<1$ then $a<g(x)<b$ as $b-a>0$, so $g:(0,1)\to(a,b)$ is well-defined. Define $k:(a,b)\to(0,1)$ by $k(y)=\frac{y-a}{b-a}$; for $a<y<b$ this lies in $(0,1)$. Then $k(g(x))=\frac{(b-a)x}{b-a}=x$ and $g(k(y))=a+(y-a)=y$. So $g$ is a bijection with inverse $k$.

\hsection{Chapter 4}
\sol{S4.1} Base $n=0$: empty sum $0=1!-1$. Step: $\sum_{k=1}^{n+1}k\cdot k!=(n+1)!-1+(n+1)(n+1)!=(n+1)!\,(n+2)-1=(n+2)!-1$.

\sol{S4.2} Base $n=0$: $6\mid0$. Step: $(n+1)^3-(n+1)=(n^3-n)+3n^2+3n=(n^3-n)+3n(n+1)$. By the IH $n^3-n=6k$; and $n(n+1)$ is even (the fact that $n(n+1)$ is even (one of $n$, $n+1$ is even)), say $2m$, so $3n(n+1)=6m$. Hence $(n+1)^3-(n+1)=6(k+m)$.

\sol{S4.3} For $0\le k<n$: $\binom nk+\binom n{k+1}=\frac{n!}{k!(n-k)!}+\frac{n!}{(k+1)!(n-k-1)!}=\frac{n!\,[(k+1)+(n-k)]}{(k+1)!(n-k)!}=\frac{(n+1)!}{(k+1)!(n-k)!}=\binom{n+1}{k+1}$. For $k=n$: $\binom nn+\binom n{n+1}=1+0=\binom{n+1}{n+1}$; for $k>n$ both sides are $0$. Induction on $n$ for $p(n)$: ``$\binom nk\in\N$ for all $k\in\N$''. $p(0)$: $\binom00=1$ and $\binom0k=0$ for $k\ge1$. Step: $\binom{n+1}0=1$, and for $k\ge0$, $\binom{n+1}{k+1}=\binom nk+\binom n{k+1}\in\N$ by the IH.

\sol{S4.4} Let $p(n)$ be the claim for $2^n\times2^n$ boards. $p(1)$: a $2\times2$ board minus one square is exactly one L-piece. Step: given a $2^{n+1}\times2^{n+1}$ board with one square removed, cut it into four $2^n\times2^n$ quadrants. The removed square lies in one quadrant. Place one L-piece at the centre of the board covering one square from each of the other three quadrants. Now every quadrant is a $2^n\times2^n$ board with exactly one square missing; tile each by $p(n)$.

\sol{S4.5} Let $p(n)$: no $d\ge2$ divides both $f_n$ and $f_{n+1}$. $p(0)$: $f_1=1$ has no divisor $\ge2$. Step: suppose $d\ge2$ divides $f_{n+1}$ and $f_{n+2}$. Then $d\mid f_{n+2}-f_{n+1}=f_n$, so $d$ divides both $f_n$ and $f_{n+1}$, contradicting $p(n)$. Hence $p(n+1)$.

\sol{S4.6} Base $n=1$: $1\ge1$. Step: by the IH the sum up to $n+1$ is at least $\sqrt n+\frac1{\sqrt{n+1}}$, so it suffices that $\sqrt n+\frac1{\sqrt{n+1}}\ge\sqrt{n+1}$. Multiplying by $\sqrt{n+1}>0$ this is $\sqrt{n(n+1)}+1\ge n+1$, i.e.\ $\sqrt{n(n+1)}\ge n$, which holds since $n(n+1)\ge n^2$ and both sides are non-negative.

\sol{S4.7} Strong induction on $n\ge1$. Since $2^0\le n$ and $2^n>n$, there is a largest $k\in\N$ with $2^k\le n$. Then $n-2^k<2^k$ (otherwise $2^{k+1}\le n$). If $n-2^k=0$, then $n=2^k$. Otherwise $1\le n-2^k<n$, so by the IH $n-2^k$ is a sum of distinct powers of $2$, each at most $n-2^k<2^k$; adding $2^k$ keeps the powers distinct.

\sol{S4.8} Base $n=1$: $f_1=1=f_3-1$. Step: $\sum_{k=1}^{n+1}f_k=(f_{n+2}-1)+f_{n+1}=f_{n+3}-1$.

\sol{S4.9} Base cases $8=3+5$, $9=3\cdot3$, $10=5\cdot2$. Step: let $n\ge11$ and assume the claim for all $8\le k<n$. Then $8\le n-3<n$, so $n-3=3a+5b$ and $n=3(a+1)+5b$. Three base cases are needed because the step reaches back by $3$.

\hsection{Chapter 5}
\sol{S5.1} Reflexive: $a=1\cdot a$. Transitive: as for linear combinations of multiples: $b=ka$, $c=lb$ give $c=(lk)a$. Not symmetric: $1\mid2$ but $2\nmid1$. Not antisymmetric: $2\mid-2$ and $-2\mid2$ but $2\ne-2$.

\sol{S5.2} Let $a\in X$; fix $b$ with $aRb$. By symmetry $bRa$. By transitivity applied to $aRb$ and $bRa$, $aRa$.

\sol{S5.3} Reflexive: $x-x=0\in\Q$. Symmetric: $y-x=-(x-y)\in\Q$. Transitive: $x-z=(x-y)+(y-z)\in\Q$. $[0]=\{y\mid 0-y\in\Q\}=\Q$. $[\sqrt2]=\{y\mid\sqrt2-y\in\Q\}=\{\sqrt2-r\mid r\in\Q\}=\{\sqrt2+r\mid r\in\Q\}$. If $[\sqrt2]=[\sqrt3]$ then $r=\sqrt3-\sqrt2\in\Q$, and $r\ne0$; then $\sqrt3+\sqrt2=1/r\in\Q$, contradicting the Chapter 0 book-level example. So $[\sqrt2]\ne[\sqrt3]$.

\sol{S5.4} Equivalence relations correspond to partitions. Partitions of a $4$-element set by block sizes: one block of $4$: $1$; sizes $3+1$: $4$ (choose the singleton); $2+2$: $3$ (choose the partner of element $1$); $2+1+1$: $6$ (choose the pair); $1+1+1+1$: $1$. Total $15$.

\sol{S5.5} Reflexive: $U\symdiff U=\varnothing$ is finite. Symmetric: $U\symdiff V=V\symdiff U$. Transitive: suppose $U\symdiff V$ and $V\symdiff W$ are finite. If $a\in U\symdiff W$, say $a\in U$, $a\notin W$ (the other case is symmetric): if $a\in V$ then $a\in V\symdiff W$; if $a\notin V$ then $a\in U\symdiff V$. So $U\symdiff W\subseteq(U\symdiff V)\cup(V\symdiff W)$, a subset of a union of two finite sets, hence finite.

\sol{S5.6} $g[X/{\sim}]=\{g([a])\mid a\in X\}=\{f(a)\mid a\in X\}=f[X]$. For injectivity: $g$ injective iff for all $a,b$, $g([a])=g([b])\Rightarrow[a]=[b]$, i.e.\ $f(a)=f(b)\Rightarrow a\sim b$, using $[a]=[b]\Leftrightarrow a\sim b$.

\sol{S5.7} ($\Leftarrow$) Suppose $m\mid n$ and $[a]_n=[a']_n$. Then $n\mid a-a'$, and $m\mid n$, so $m\mid a-a'$ by transitivity of $\mid$; thus $[a]_m=[a']_m$. ($\Rightarrow$) If the rule is well-defined, then since $[0]_n=[n]_n$ we must have $[0]_m=[n]_m$, i.e.\ $m\mid n-0=n$.

\sol{S5.8} ($\Rightarrow$) An equivalence relation is reflexive; and if $aRb$ and $aRc$, then $bRa$ by symmetry, and $bRa$, $aRc$ give $bRc$ by transitivity. ($\Leftarrow$) Suppose $R$ is reflexive and has the stated property. Symmetric: let $aRb$. Also $aRa$. Applying the property to $aRb$ and $aRa$ gives $bRa$. Transitive: let $aRb$ and $bRc$. By symmetry $bRa$. Applying the property to $bRa$ and $bRc$ gives $aRc$.
